% Compile twice with pdflatex psd_product.tex to resolve the references.
\documentclass[10pt,twoside,reqno]{amsart}
\usepackage[T1]{fontenc}
\usepackage{lmodern}
\usepackage{amsmath,amssymb,amsthm}
\usepackage[letterpaper,textwidth=6.2in,textheight=9in,centering]{geometry}

\newtheorem{theorem}{Theorem}
\newtheorem{proposition}{Proposition}[section]
\newtheorem{lemma}[proposition]{Lemma}
\numberwithin{equation}{section}
\newcommand{\C}{\mathbb C}
\newcommand{\R}{\mathbb R}
\newcommand{\T}{\mathbb T}
\newcommand{\Real}{\operatorname{Re}}
\newcommand{\dd}{\,\mathrm{d}}
\setlength{\emergencystretch}{1em}

\title[A product inequality for positive semidefinite matrices]
      {A product inequality for positive semidefinite matrices}
\author[Saigyukari]{Saigyukari}
\date{}
\pagestyle{headings}

\begin{document}
\begin{abstract}
For every complex positive semidefinite matrix of order \(n\), we bound the
product of the distances of its entries from \(1\) by the corresponding
product for the \(n\)th powers of its diagonal entries. The bound holds in
every rank and is attained by a rank-one matrix of roots of unity.
A radial deformation reduces the inequality to a Poisson-kernel estimate
whose harmonic minorant has explicit rank-one Fourier coefficients in
logarithmic radii. A sum-of-squares factorization reduces positivity in every
order to the same fourth-order differential inequality. We prove this inequality analytically by a log-concave density comparison
and a scalar differential equation. The complete proof is formalized in
Lean~4 with Mathlib and uses only Lean's standard axioms.
\end{abstract}
\maketitle

\section{Introduction and diagonal reduction}

\begin{theorem}\label{thm:product}
For every \(n\ge1\) and every positive semidefinite \(A=(a_{ij})\in\C^{n\times n}\),
\begin{equation}\label{eq:product}
 \prod_{i,j=1}^{n}|1-a_{ij}|
 \ \le\ \prod_{i=1}^{n}|1-a_{ii}^{\,n}|.
\end{equation}
\end{theorem}

Equality holds for \(a_{ij}=d\,\omega^{i-j}\), where \(d\ge0\) and
\(\omega\) is a primitive \(n\)th root of unity. Indeed, \(A=d\,u u^*\)
with \(u_i=\omega^i\), and \(\prod_j(1-d\,\omega^{i-j})=1-d^n\) in each row.

Write \(d_i=a_{ii}\) and \(S_m(z)=\sum_{\ell=0}^{m-1}z^\ell\).
If any \(d_i=1\), both sides of \eqref{eq:product} vanish. Otherwise division
by \(\prod_i|1-d_i|\) reduces it to
\begin{equation}\label{eq:cancel}
 \prod_{i\ne j}|1-a_{ij}|\ \le\ \prod_i S_n(d_i).
\end{equation}
The latter inequality extends to unit diagonal entries by continuity, using
\(\rho A\to A\) through values with \(\rho d_i\ne1\) for every \(i\).
Order one is an identity. In order two, positivity gives
\(|a_{12}|^2\le d_1d_2\), whence
\[
 |1-a_{12}|^2\le(1+\sqrt{d_1d_2})^2\le(1+d_1)(1+d_2).
\]
The difference in the last inequality is \((\sqrt{d_1}-\sqrt{d_2})^2\).

For \(n\ge3\) put \(t_i=\max(1,d_i)\) and \(b_{ij}=a_{ij}/(t_it_j)\).
The matrix \(B=(b_{ij})\) is positive semidefinite by diagonal congruence.
For \(T=t_i t_j\), the bound \(|a_{ij}|^2\le d_i d_j\le T\) gives
\[
 T|1-b_{ij}|^2-|1-a_{ij}|^2=(T-1)(1-|a_{ij}|^2/T)\ge0.
\]
If \(d_i>1\), then \(b_{ii}=1/d_i\), so reversing the geometric sum gives
\(S_n(d_i)=t_i^{n-1}S_n(b_{ii})\); the same identity holds when \(d_i\le1\).
Multiplication over the ordered pairs \(i\ne j\) contributes
\(\prod_{i\ne j}\sqrt{t_i t_j}=\prod_i t_i^{\,n-1}\). Thus
\eqref{eq:cancel} for \(B\) implies it for \(A\).
Since \(0\le b_{ii}\le1\), applying the bound first to \(\rho B\),
\(0<\rho<1\), and letting \(\rho\uparrow1\) reduces the proof to
\(0\le a_{ii}<1\).

\section{Radial deformation and harmonic minorants}

Write \(x_i=a_{ii}<1\) and \(a_{ij}=\sqrt{x_i x_j}\,c_{ij}\), where
\(c_{ij}=\langle u_i,u_j\rangle\) for unit vectors \(u_i\). This follows
from a Gram representation of \(A\); at a zero Gram vector any unit direction
will do. Let \(P(\zeta)=\Real((1+\zeta)/(1-\zeta))\) for \(|\zeta|<1\).
Hold the \(u_i\) fixed and evolve \(\dot x_i=x_i w(x_i)\), where \(w\) is
positive and continuous on \([0,1)\) and locally Lipschitz there. With the
prescribed radii at time zero, a positive coordinate satisfies
\[
 -t=\int_{x_i(t)}^{x_i(0)}\frac{\dd u}{u w(u)}.
\]
The integral diverges at zero because \(w(0)>0\). Thus the backward flow
exists for all \(t\le0\) and tends to zero. A zero coordinate stays zero.

Set \(w_i=w(x_i)\) and
\(Q=\sum_i\log(1-x_i^n)-\sum_{i,j}\log|1-a_{ij}|\).
Along the flow \(\dot a_{ij}=(w_i+w_j)a_{ij}/2\).
Using \(2\Real(\zeta/(1-\zeta))=P(\zeta)-1\) and
\(\sum_{i,j}(w_i+w_j)/2=n\sum_iw_i\), differentiation gives
\begin{equation}\label{eq:flow}
 2\dot Q=\sum_{i,j}\frac{w_i+w_j}{2}P(a_{ij})
               -n\sum_i w_i P(x_i^n).
\end{equation}
A harmonic minorant can control the two sums at the same time.

\begin{proposition}\label{prop:flow}
Let \(\eta\) be real-valued on \([0,1)\), and let \(v_k(x)\) take values in
real Hilbert spaces, \(1\le k\le n-2\). Suppose
\[
 h(x,y,z)=\frac{\eta(x)+\eta(y)}2+
       \sum_{k=1}^{n-2}\langle v_k(x),v_k(y)\rangle\,\Real(z^k)
 \ \le\ \frac{w(x)+w(y)}2\,P(\sqrt{xy}\,z)
\]
for \(0\le x,y<1\), \(z\in\T=\{z:|z|=1\}\), with equality if \(x=y\)
and \(z^n=1\), \(z\ne1\). Then \eqref{eq:product} holds in order \(n\).
\end{proposition}

\begin{proof}
Both sides are harmonic on the disk and continuous on its closure.
The minimum principle extends the bound to \(|z|\le1\). For every \(k\),
complexifying the real coefficient space in the tensor product gives
\[
 \sum_{i,j}\langle v_k(x_i),v_k(x_j)\rangle\,\Real(c_{ij}^k)
 =\Big\|\sum_i v_k(x_i)\otimes u_i^{\otimes k}\Big\|^2\ge0.
\]
Consequently \(\sum_{i,j}h(x_i,x_j,c_{ij})\ge n\sum_i\eta(x_i)\).
To keep the exact diagonal contribution, let \(\mu_n=\{z:z^n=1\}\).
Summing \(P(xz)=1+2\Real\sum_{\ell\ge1}x^\ell z^\ell\) over \(\mu_n\)
retains precisely the multiples of \(n\), so \(\sum_{z\in\mu_n}P(xz)=nP(x^n)\).
Likewise \(\sum_{z\in\mu_n}h(x,x,z)=n\eta(x)\). All the differences except
the one at \(z=1\) vanish by hypothesis. Hence
\[
 w(x)P(x)-h(x,x,1)=n\bigl(w(x)P(x^n)-\eta(x)\bigr).
\]
Adding these diagonal differences to the lower bound for the sum of \(h\)
shows that the first sum in \eqref{eq:flow} is at least \(n\sum_i w_iP(x_i^n)\).
Thus \(\dot Q\ge0\). Since \(Q\to0\) along the backward flow,
\(Q(0)\ge0\), which is \eqref{eq:product} for the reduced matrix.
Section~1 gives the result for all diagonal entries.
\end{proof}

We use one scalar for each nonconstant Fourier coefficient. For \(m\ge2\) set
\begin{equation}\label{eq:weight}
 \begin{split}
 D_m(x)&=1-x^{2m}-m(1-x^2)x^{m-1},\qquad
 w_n(x)=\frac{(1-x^n)^2}{D_n(x)},\\
 p_k(x)&=x^{k/2}\sqrt{\frac{D_{n-k}(x)}{D_n(x)}},\qquad
 h_n(x,y,z)=1+2\sum_{k=1}^{n-2}p_k(x)p_k(y)\,\Real(z^k).
 \end{split}
\end{equation}
Here \(w_n(0)=1\) and \(p_k(0)=0\). We shall prove on \(\T\) the stronger
geometric-mean inequality
\begin{equation}\label{eq:minorant}
 h_n(x,y,z)\le\sqrt{w_n(x)w_n(y)}\,P(\sqrt{xy}\,z).
\end{equation}

\section{Logarithmic radii and a sum of squares}

For \(0<x,y<1\) put \(a=-\log x\), \(b=-\log y\), \(c=(a+b)/2\),
\(d=(a-b)/2\). For \(q>0\) and real \(s\ge1\) write
\[
 H_s(q)=\sinh(sq)-s\sinh q,\qquad
 W_n(q)=\frac{\cosh(nq)-1}{H_n(q)},\qquad
 U_s=\sqrt{H_s(a)H_s(b)},\qquad U_0=0 .
\]
We have \(U_1=0\) and \(H_s(q)>0\) for \(s>1\), since \(\sinh q/q\)
is strictly increasing for \(q>0\). The identities
\[
 D_m(e^{-q})=2e^{-mq}H_m(q),\qquad w_n(e^{-q})=W_n(q)
\]
also establish the positivity of the denominators in \eqref{eq:weight}.
In particular \(w_n\) is smooth on \([0,1)\).
For \(z=e^{i\theta}\), formula \eqref{eq:weight} and the Poisson kernel become
\[
 h_n=1+2\sum_{k=1}^{n-2}\frac{U_{n-k}}{U_n}\cos(k\theta),\qquad
 P(e^{-c}z)=\frac{\sinh c}{\cosh c-\cos\theta}.
\]

Multiplication by the Poisson denominator expresses the remainder in squares
of the geometric sums from Section~1. Put \(L_0=0\) and define
\begin{align}
 L_m&=\tfrac12(U_{m+1}+U_{m-1})-\cosh c\,U_m \quad (m\ge1), \notag\\
 B_1&=L_1,\qquad B_j=L_j-2L_{j-1}+L_{j-2}\quad (j\ge2). \label{eq:weights}
\end{align}
For a real number \(M\) let
\[
 R_M=M P(e^{-c}z)-h_n,\qquad
 E_n=M\sinh c\,U_n-\tfrac12(U_{n+1}-U_{n-1}).
\]

\begin{proposition}\label{prop:squares}
For every \(z=e^{i\theta}\) and every \(M\),
\begin{equation}\tag{A}\label{eq:squares}
 (\cosh c-\cos\theta)U_n R_M
   =E_n+\sum_{j=1}^n B_j\,|S_{n-j+1}(z)|^2 .
\end{equation}
\end{proposition}

\begin{proof}
In \(U_n h_n\) the coefficient of \(z^k\) is \(U_{n-|k|}\) for
\(|k|\le n-2\), and zero otherwise. Multiplying by
\(\cos\theta-\cosh c=(z+z^{-1})/2-\cosh c\) gives the coefficient
\(L_{n-|k|}\) for \(1\le|k|\le n-1\), including the endpoint because
\(U_0=U_1=0\). The constant becomes \(U_{n-1}-\cosh c\,U_n\).
Adding \(M\sinh c\,U_n\) changes this constant to \(E_n+L_n\).

By \eqref{eq:weights}, \(\sum_{j=1}^m B_j=L_m-L_{m-1}\).
Summing once more gives
\(L_m=\sum_{j=1}^m(m-j+1)B_j\). Since
\(|S_\ell(z)|^2=\sum_{|k|<\ell}(\ell-|k|)z^k\),
the coefficient of \(z^k\) in the sum of squares is
\(\sum_{j=1}^{n-|k|}(n-|k|-j+1)B_j=L_{n-|k|}\) for \(|k|<n\).
This matches every coefficient, including the constant after adding \(E_n\).
\end{proof}

\section{Positivity of the coefficients}

The constant remainder in \eqref{eq:squares} is controlled by the
variation of one ratio of hyperbolic functions.

\begin{lemma}\label{lem:remainder}
For \(M=\sqrt{W_n(a)W_n(b)}\), one has \(E_n\ge0\).
\end{lemma}

\begin{proof}
Put \(r(q)=\tfrac12\log(H_{n+1}(q)/H_{n-1}(q))\). We first prove
\begin{equation}\label{eq:lipschitz}
 r(q)\ge q,\qquad 0\le r'(q)\le1 .
\end{equation}
Write \(u=n+1\), \(v=n-1>1\). The coefficients of \(q^{2\ell+1}\) in
\(H_u,H_v\) are \(\alpha_\ell=(u^{2\ell+1}-u)/(2\ell+1)!\) and
\(\beta_\ell=(v^{2\ell+1}-v)/(2\ell+1)!\), respectively, for \(\ell\ge1\).
Up to the positive factor \(u(u^2-1)/(v(v^2-1))\), their ratio is
\(\sum_{j<\ell}u^{2j}/\sum_{j<\ell}v^{2j}\), with \(j\) starting at zero.
This weighted mean of the increasing numbers \((u/v)^{2j}\) increases with
\(\ell\). Pairing indices \(\ell>k\) in \(H_u'H_v-H_uH_v'\) gives the
terms \(2(\ell-k)(\alpha_\ell\beta_k-\alpha_k\beta_\ell)q^{2\ell+2k+1}\ge0\).
The series converge absolutely; hence \(r'\ge0\).

Let \(h_k(X)\) denote the complete homogeneous symmetric polynomial,
defined by \(\sum_{k\ge0}h_k(X)y^k=\prod_i(1-X_i y)^{-1}\).
Summing the two terms of \(H_m(q)\) separately gives
\[
 \sum_{m\ge0} H_m(q)y^m
 =y\sinh q\left(\frac1{1-2y\cosh q+y^2}-\frac1{(1-y)^2}\right)
 =\frac{H_2(q)y^2}{(1-e^q y)(1-y)^2(1-e^{-q}y)},
\]
where in this series \(H_0(q)=0\) and \(H_2(q)=2(\cosh q-1)\sinh q\).
Thus \(H_m(q)=H_2(q)h_{m-2}(e^q,1,1,e^{-q})\) for \(m\ge2\).
Homogeneity now yields
\begin{equation}\label{eq:homogeneous}
 e^{2(r(q)-q)}
 =\frac{h_{n-1}(X)}{h_{n-3}(X)},\qquad
 X=(1,e^{-q},e^{-q},e^{-2q}).
\end{equation}

For completeness, the \(h_k\) form a log-concave sequence for positive
variables. This is an equality for one variable. Adjoining a variable \(z\)
to a positive log-concave sequence \(f_k\) of complete homogeneous polynomials
gives \(g_k=\sum_{j=0}^k z^j f_{k-j}=f_k+zg_{k-1}\). Therefore
\[
 g_k^2-g_{k-1}g_{k+1}
 =g_k f_k-g_{k-1}f_{k+1}
 =z^k f_0f_k+\sum_{j=0}^{k-1}z^j(f_{k-j}f_k-f_{k-1-j}f_{k+1})\ge0,
\]
since the successive ratios \(f_i/f_{i-1}\) decrease.
For any one coordinate \(z\), differentiation of the generating function gives
\(\partial_z\log h_k=\sum_{j=1}^k z^{j-1}h_{k-j}/h_k\).
Log-concavity makes each ratio \(h_{k-j}/h_k\) nondecreasing in \(k\).
Thus \(\partial_z\log(h_{k+1}/h_k)\ge0\), proving coordinatewise
monotonicity, also at nonnegative coordinates wherever the denominator is
positive, by continuity.

The quotient in \eqref{eq:homogeneous} is a product of two consecutive such
ratios. Its coordinates decrease or stay fixed with \(q\), and its limit at
\(q=\infty\) is \(1\). Hence \(r(q)-q\) decreases to zero, giving the
remaining bounds in \eqref{eq:lipschitz}.

Finally put \(T(q)=\sqrt{H_{n+1}(q)H_{n-1}(q)}/H_n(q)\). Subtraction gives
\[
 H_{n+1}(q)-H_{n-1}(q)
 =2\sinh q\,\cosh(nq)-2\sinh q,\qquad
 T(q)\sinh r(q)=W_n(q)\sinh q.
\]
For \(u,v>0\) the identity
\[
 \frac{\sinh^2((u+v)/2)}{\sinh u\,\sinh v}
 =1+\frac{\sinh^2((u-v)/2)}{\sinh u\,\sinh v}
\]
follows from the hyperbolic addition formulas. Its right side at
\((r(a),r(b))\) is at most its value at \((a,b)\), since
\(|r(a)-r(b)|\le|a-b|\), \(r(a)\ge a\), and \(r(b)\ge b\).
Taking square roots and using the identity for \(T\) gives
\[
 \frac{U_{n+1}-U_{n-1}}{2U_n}
 =\sqrt{T(a)T(b)}\,\sinh\!\left(\frac{r(a)+r(b)}2\right)
 \le\sqrt{W_n(a)W_n(b)}\,\sinh c .
\]
This is the asserted sign of \(E_n\).
\end{proof}

The first three coefficients involve the endpoint \(U_0=U_1=0\) and
can be computed directly.

\begin{lemma}\label{lem:first}
With \(D=\cosh d\), one has
\(B_1=U_2/2\), \(B_2=U_2(D-1)\), and \(B_3\ge2U_2D(D-1)\).
\end{lemma}

\begin{proof}
Write \(C=\cosh c\) and \(A_q=\cosh q\). Since
\(H_2(q)=2\sinh q(A_q-1)\) and \(H_3(q)=4\sinh^3 q\),
\[
 \frac{U_3}{U_2}
 =2\sqrt{(1+A_a)(1+A_b)}
 =4\cosh(a/2)\cosh(b/2)=2(C+D).
\]
Here \(L_1=U_2/2\) and \(L_2=U_3/2-CU_2\); subtracting \(2L_1\) from
\(L_2\) gives the formula for \(B_2\).
Also \(H_4(q)=4\sinh q(2A_q^3-A_q-1)\), so
\[
 \frac{H_4(q)}{2H_2(q)}
 =2A_q^2+2A_q+1=A_q^2+(1+A_q)^2 .
\]
Cauchy--Schwarz applied to \((A_a,1+A_a)\) and \((A_b,1+A_b)\) yields
\[
 \frac{U_4}{2U_2}\ge A_aA_b+(1+A_a)(1+A_b)
 =2C^2+2D^2+2CD-1,
\]
using \(A_a+A_b=2CD\) and \(A_aA_b=C^2+D^2-1\).
Expanding \(B_3=L_3-2L_2+L_1\) and substituting the two ratios gives
\begin{align*}
 \frac{B_3}{U_2}
 &=\frac{U_4}{2U_2}-(C+1)\frac{U_3}{U_2}+2C+1\\
 &\ge 2C^2+2D^2+2CD-1-2(C+1)(C+D)+2C+1=2D(D-1).
\end{align*}
\end{proof}

For the remaining coefficients, two centered differences give an integral
with a nonnegative kernel.

\begin{lemma}\label{lem:integral}
Put \(U(s)=U_s\) for \(s\ge1\), \(K_0(u)=1-|u|\), and
\(K_c(u)=\sinh(c(1-|u|))/c\) on \([-1,1]\). For every integer \(m\ge3\),
\begin{equation}\label{eq:integral}
 2B_{m+1}=\int_{-1}^1\int_{-1}^1 K_0(u)K_c(v)
       (U^{(4)}-c^2U'')(m+u+v)\dd v\dd u .
\end{equation}
\end{lemma}

\begin{proof}
The analytic factorization \(H_s(q)=(s-1)\widetilde H_q(s)\) has
\(\widetilde H_q(1)=q\cosh q-\sinh q=\int_0^q u\sinh u\dd u>0\).
Thus \(U(s)=(s-1)\sqrt{\widetilde H_a(s)\widetilde H_b(s)}\) on \(s\ge1\)
extends real-analytically across \(1\), including the corner
\(m=3,\ u=v=-1\) in \eqref{eq:integral}.

For \(h=c\) or \(0\), define
\(\Delta_h F(s)=F(s+1)+F(s-1)-2\cosh h\,F(s)\).
On each half of \([-1,1]\), \(K_h''=h^2K_h\). Its endpoint values vanish,
its endpoint derivatives are \(1,-1\), and
\(K_h'(0+)-K_h'(0-)=-2\cosh h\). Put
\(G(u)=K_h(u)F'(s+u)-K_h'(u)F(s+u)\). Integration by parts gives
\[
 \int_{-1}^1 K_h(u)(F''-h^2F)(s+u)\dd u
 =[G]_{-1}^{0-}+[G]_{0+}^{1}=\Delta_h F(s).
\]
Now \(2L_m=\Delta_c U(m)\), and \(2B_{m+1}=\Delta_0\Delta_c U(m)\).
Applying the integral formula twice gives \eqref{eq:integral}.
\end{proof}

The differential expression in \eqref{eq:integral} is independent of the
matrix order. Its sign follows from two comparisons of logarithmic
derivatives, one using the moments of a density and the other using a
scalar differential equation.

\begin{proposition}\label{prop:differential}
Let \(a,b>0\), \(s>1\), \(c=(a+b)/2\), and \(d=(a-b)/2\).
For \(F_q(s)=H_s(q)\), put
\[
 w_q=\frac{q(\sinh q-qe^{-sq})}{F_q(s)},\qquad
 \overline w=\frac{w_a+w_b}{2}.
\]
Then \(U=\sqrt{F_aF_b}\) satisfies
\begin{equation}\label{eq:diff-bound}
 U^{(4)}-c^2 U''\ge \frac{d^2}{3}\,\overline w\,U\ge0.
\end{equation}
Equality in \(U^{(4)}-c^2U''\ge0\) holds exactly when \(a=b\).
\end{proposition}

\begin{proof}
By symmetry assume \(a\ge b\). Primes on \(F_q\) and its logarithmic
derivatives refer to \(s\). A bar denotes the average of the values at
\(a,b\), and \(\Delta\) denotes half their difference.

\medskip\noindent
\textit{1. Fourth derivative.}
Set
\[
 p_q=\frac{F_q'}{F_q},\qquad A_q=\frac{F_q''}{F_q},\qquad
 B_q=\frac{F_q'''}{F_q},\qquad C_q=q^2 A_q.
\]
Write \(P=\overline p\), \(D=\Delta p\), \(S=P-c\), and
\(j_q=q^3e^{-sq}/F_q\).
Since \(\cosh(sq)-\sinh(sq)=e^{-sq}\),
\[
 A_q=qp_q+w_q,\qquad B_q=q A_q+j_q,\qquad
 p_q'=A_q-p_q^2.
\]
Here \(w_q>0\) since \(\sinh q\ge q>qe^{-sq}\).
The function \(F'_q-qF_q\) has derivative \(F_qw_q>0\) on \((1,\infty)\)
and has value \(q\cosh q-\sinh q>0\) at \(1\).
Hence \(p_q>q\) and \(S>0\).

The logarithmic derivatives of \(F_a/U\) and \(F_b/U\) are \(D\) and \(-D\).
The product rule in \(F_q=U(F_q/U)\), divided by \(F_q\) and averaged or
subtracted, gives
\[
 \begin{aligned}
 U''/U&=\overline A-D^2,\\
 \Delta B&=3(U''/U)D+3P D'+D''+D^3,\\
 \overline C&=U^{(4)}/U+6(U''/U)D^2+12PDD'
                +4DD''+3(D')^2+D^4.
 \end{aligned}
\]
Subtracting \(c^2U''/U\) and eliminating \(D''\) leaves
\[
 \frac{U^{(4)}-c^2U''}{U}
 =\overline C-c^2\overline A-4D\Delta B
     +(6\overline A+c^2)D^2-3D^4-3(D')^2.
\]
Use \(\overline C-c^2\overline A=d^2\overline A+2cd\Delta A\),
\(\Delta B=c\Delta A+d\overline A+\Delta j\), and
\(\Delta A=D'+2PD\). This gives
\[
 (d^2-4dD+6D^2)\overline A+2c(d-2D)(D'+2PD)
     +c^2D^2-4D\Delta j-3D^4-3(D')^2.
\]
Now \(\overline A=c^2+cS+dD+\overline w\). With
\(Z=c(d-D)-D'=(2D-d)S-\Delta w\), the expression becomes
\begin{equation}\label{eq:diff-factor}
 \begin{split}
 \frac{U^{(4)}-c^2U''}{U}
 &=\overline w\bigl((d-2D)^2+2D^2\bigr)-4D\Delta j\\
 &\quad +(d-D)D(d^2-3dD+3D^2)\\
 &\quad +Z\bigl(c(d+D)+3D'\bigr)+cS(d^2-2D^2).
 \end{split}
\end{equation}
For the last regrouping, the terms in \(c^2,D'\) are
\(c^2(d^2-D^2)+2c(d-2D)D'-3(D')^2\);
substituting \(Z=c(d-D)-D'\) in \eqref{eq:diff-factor} and using
\(Z=(2D-d)S-\Delta w\) gives the displayed form after collection.

We will prove
\begin{equation}\label{eq:diff-comparisons}
 0\le D\le d,\qquad 0\le D'\le cd/2,\qquad
 \Delta j\le0,\quad \Delta w\le0 .
\end{equation}
Then the second and third terms in \eqref{eq:diff-factor} are nonnegative,
since \(d^2-3dD+3D^2=(d-3D/2)^2+3D^2/4\).
If \(D\le d/2\), then \(Z\ge c(d/2-D)\ge0\) and
\(d^2-2D^2\ge0\). If \(D>d/2\), then \(Z\ge(2D-d)S\ge0\).
The last two terms together are at least
\[
 cS\bigl((2D-d)(d+D)+d^2-2D^2\bigr)=cS dD\ge0.
\]
Finally \((d-2D)^2+2D^2=6(D-d/3)^2+d^2/3\), so
\eqref{eq:diff-comparisons} implies \eqref{eq:diff-bound}.

\medskip\noindent
\textit{2. A density comparison.}
For \(q>0\) put
\[
 f_q(v)=e^{-sv}(\cosh(\min(v,2q-v))-1)\quad (0\le v\le 2q),
 \qquad M_q=\int_0^{2q} f_q(v)\dd v.
\]
The density is positive on the open interval. Write \(\mu_q\) and
\(\mathcal V_q\) for the mean and variance with respect to \(f_q/M_q\).
We have the representation
\begin{equation}\label{eq:diff-density}
 F_q(s)=\tfrac12 s(s^2-1)e^{sq}M_q.
\end{equation}
Indeed \(K(u)=\cosh(q-|u|)-1\) on \([-q,q]\) satisfies \(K''-K=1\)
on each half. Its value and derivative vanish at the outer endpoints,
and \(K'(0+)-K'(0-)=-2\sinh q\).
Integration by parts on the two halves gives
\(\int_{-q}^q e^{su}\dd u=(s^2-1)\int_{-q}^q e^{su}K(u)\dd u
 +2\sinh q\). Substituting \(v=q-u\) proves \eqref{eq:diff-density}.
Differentiation in \(s\) under the finite integral yields
\[
 p_q=q+\ell'-\mu_q,\qquad p_q'=\ell''+\mathcal V_q,qquad
 \ell(s)=\log(s(s^2-1)).
\]

Both moments are nondecreasing in \(q\). To see this, write
\(g_q=f_q/M_q\), \(\lambda=\partial_q M_q/M_q\ge0\).
The left half of \(f_q\) is independent of \(q\); on the right half
\(\partial_q f_q=2e^{-sv}\sinh(2q-v)=-2\partial_v f_q-2s f_q\).
The boundary terms at \(v=q\) cancel when differentiating \(M_q\), and
\(f_q(2q)=0\), so
\(\partial_q M_q=2f_q(q)-2s\int_q^{2q}f_q\).
Define
\[
 J_q(v)=
 \begin{cases}
  \lambda\displaystyle\frac{\int_0^v f_q(u)\dd u}{f_q(v)},&0<v\le q,\\[6pt]
  2-(2s+\lambda)\displaystyle\frac{\int_v^{2q} f_q(u)\dd u}{f_q(v)},
    &q\le v<2q .
 \end{cases}
\]
The formula for \(\partial_q M_q\) gives agreement at \(q\).
Direct differentiation gives \(\partial_q g_q=-\partial_v(J_q g_q)\)
on the two halves, with vanishing flux at the outer endpoints.

The function \(f_q\) is log-concave: \(\log(\cosh t-1)\) is increasing
and concave for \(t>0\), since it equals \(\log 2+2\log\sinh(t/2)\).
Compose with the concave tent \(\min(v,2q-v)\) and add \(-sv\).
For a log-concave positive density \(f\) on an interval, the ratio
\(f(v-h)/f(v)\) is nondecreasing in \(v\) for \(h\ge0\) on its domain,
by concavity of \(\log f\). Integrating in \(h\), with the integration
length increasing in \(v\), shows that the left integral-to-density
ratio is nondecreasing. The same argument with \(f(v+h)/f(v)\) gives
a nonincreasing right ratio. Thus \(J_q\) is nonnegative and
nondecreasing. Splitting the integrals at \(q\), differentiating in
\(q\), and integrating the flux by parts now gives
\[
 \partial_q\mu_q=\int_0^{2q} J_q g_q\dd v\ge0,\qquad
 \partial_q\mathcal V_q=2\int_0^{2q}(v-\mu_q)
           (J_q(v)-J_q(\mu_q))g_q(v)\dd v\ge0 .
\]
The subtraction of \(J_q(\mu_q)\) uses \(\int(v-\mu_q)g_q=0\).
It follows that \(D=d-(\mu_a-\mu_b)/2\le d\) and
\(D'=(\mathcal V_a-\mathcal V_b)/2\ge0\).

\medskip\noindent
\textit{3. Comparison of normalized integrals.}
Put \(L(x)=x\cosh x-\sinh x\), \(T(x)=L(x)/\sinh x\), and
\(k(x)=x^2/T(x)\) for \(x>0\).
The power series for \(x^2\sinh x\) and \(L(x)\) have coefficient
ratios \(2m+1\) in degree \(2m+1\), \(m\ge1\). Hence \(k\ge3\) and
\(k'\ge0\). We use here the elementary series comparison rule that
increasing coefficient ratios give an increasing quotient of sums
on the positive axis. For series \(f=\sum f_i x^i,\ g=\sum g_i x^i\) with
positive denominator coefficients at the degrees used, pair the terms
of \(f'g-fg'\) with indices \(i>j\) to get
\((i-j)(f_i g_j-f_j g_i)x^{i+j-1}\).
Decreasing ratios reverse the conclusion.

Logarithmic differentiation and \(L'=x\sinh x\) give
\(xk'=k(3+T-k)\).
Since
\(T(T+2)=x^2+x^2/\sinh^2 x-1\le x^2=kT\), we have \(T\le k-2\).
Thus
\begin{equation}\label{eq:diff-k}
 x\bigl((k-2)k'-x\bigr)=k(k-3)(T-k+2)\le0.
\end{equation}
The function \(\kappa(z)=k(\sqrt z)\) is concave for \(z>0\).
Indeed \(2z\kappa'=z+3\kappa-\kappa^2\), so
\(2z\kappa''=1-(2\kappa-1)\kappa'\).
Set \(R=(2\kappa-1)\kappa'-1\). It satisfies
\[
 (zR)'=2z(\kappa')^2-\frac{2\kappa-3}{2z}\,zR.
\]
If \(zR\) were negative at some \(z_0>0\), it would attain a negative
minimum on \([\varepsilon,z_0]\) away from \(\varepsilon\) for
sufficiently small \(\varepsilon>0\), since \(zR\ge-z\).
At such a minimum its left derivative is nonpositive, whereas the
displayed derivative is positive. This proves \(R\ge0\) and concavity.

For \(0<u<1\) define
\[
 \gamma_x(t)=\frac{L(tx)}{t^2 L(x)},\qquad
 m_x(u)=\left(\int_u^1\gamma_x(t)\dd t\right)^{-1}
       =\frac{L(x)}{\sinh x-\sinh(ux)/u}.
\]
The integral identity follows by differentiating \(\sinh(tx)/t\).
For \(x\ge y>0\) set \(\eta(t)=\log(\gamma_x(t)/\gamma_y(t))\).
We have \(\eta(1)=0\) and
\[
 \frac{\eta'(t)}t=\frac{k(tx)-k(ty)}{t^2}
       =\int_{y^2}^{x^2}\kappa'(t^2 z)\dd z.
\]
This is nonnegative and nonincreasing in \(t\), so \(\eta\) is increasing
and concave as a function of \(t^2\).
The mean of \(t^2\) under the probability density proportional to
\(\gamma_y(t)\) on \([u,1]\) is at least \((u^2+1)/2\).
Indeed \(\gamma_y(t)/t\) increases by \(k\ge3\) and \(L'=x\sinh x\).
Against the weight \(t\dd t\), the mean of \(t^2\) is \((u^2+1)/2\).
Multiplying the weight by an increasing function can only increase the
mean, by integrating the product of the differences of the two functions
at two points.
Concavity gives \(\eta(t)\ge\eta(u)(1-t^2)/(1-u^2)\).
Since \(\eta(u)\le0\), its mean is at least \(\eta(u)/2\).
The tangent inequality for the exponential then implies
\[
 \frac{\int_u^1\gamma_x(t)\dd t}{\int_u^1\gamma_y(t)\dd t}
 \ge e^{\eta(u)/2}=\sqrt{\gamma_x(u)/\gamma_y(u)} .
\]
Thus \(m_x'(u)=\gamma_x(u)m_x(u)^2\le m_y'(u)\).
Now \(\gamma_x(1)=1\) and \(\gamma_x'(1)=k(x)-2\). Integrating the
first-order Taylor expansion at \(1\) and taking the reciprocal gives
\[
 m_x(u)=\frac1{1-u}+\frac{k(x)-2}{2}+O(1-u),\qquad
 m_x(u)-m_y(u)\ge\frac{k(x)-k(y)}2 .
\]

Apply this to \(x=sa,\ y=sb,\ u=1/s\). Direct substitution gives
\(sp_q-1=m_{sq}(u)\) and \(s^2 A_q=k(sq)m_{sq}(u)\).
Let \(W=sD\), \(e=(k(x)-k(y))/2\), and let \(k_0,m_0\) be the averages
of \(k(x),k(y)\) and \(m_x(u),m_y(u)\), respectively.
The comparison just proved gives \(W\ge e/2\ge0\).
Also \(m_z\ge T(z)\) by the formula for \(m_z\).
The series of \(L(z)\) gives \(L(z)\ge2(\sinh z-z)\), and
\(\sinh(uz)/u\ge z\), so \(m_z\ge2\).
Since \(k'\ge0\) implies \(k-2\le T+1\), we have \(2m_0\ge k_0-2\).
Taking half-differences in \(p_q'=A_q-p_q^2\) gives
\[
 s^2D'=e m_0+(k_0-2-2m_0)W
       \le (k_0-2)e/2
       =\frac{(k(x)-2)^2-(k(y)-2)^2}{8}.
\]
By \eqref{eq:diff-k} the last quotient is at most
\((x^2-y^2)/8=s^2cd/2\). This proves \(D\ge0\) and \(D'\le cd/2\).

\medskip\noindent
\textit{4. The remainders.}
The series of \(F_q/q^3\) has nonnegative coefficients for \(s>1\).
Thus \(j_q=e^{-sq}/(F_q/q^3)\) decreases in \(q\).
For \(w_q\), use the equivalent quotient
\[
 w_q=\frac{q(e^{(s+1)q}-e^{(s-1)q}-2q)}
             {e^{2sq}-1-s(e^{(s+1)q}-e^{(s-1)q})}.
\]
Let \(x=(s+1)/(2s)\), \(y=(s-1)/(2s)\), and
\(h_m=(x^{m+1}-y^{m+1})/(x-y)\). Then \(x+y=1\), \(x-y=1/s\), and
\[
 h_0=h_1=1,\qquad h_{m+2}=h_{m+1}-xyh_m,\qquad
 h_m^2-h_{m-1}h_{m+1}=(xy)^m\quad(m\ge1).
\]
Consequently \(h_m>0\) is decreasing and log-concave.
The denominator starts in degree \(k=3\). In degree \(k\ge3\),
the numerator and denominator coefficients have ratio
\[
 \frac{k h_{k-2}}{2s^2(1-h_{k-1})}
 =\frac{k h_{k-2}}{2s^2xy\sum_{r=0}^{k-3}h_r}.
\]
Indeed the denominator coefficient is \((2s)^k(1-h_{k-1})/k!\),
and the recurrence for \(h_m\) telescopes to
\(1-h_{k-1}=xy\sum_{r=0}^{k-3}h_r\).
For \(m=k-2\ge1\), the inverse of the displayed ratio divided by
\(2s^2xy\) equals
\[
 \frac{m}{m+2}\cdot\frac{1}{m}\sum_{i=1}^m\frac{h_{m-i}}{h_m}.
\]
For fixed \(i\), the summand is nondecreasing in \(m\) by log-concavity.
The added summand when \(m\) grows is largest by monotonicity of \(h_m\).
Hence the average increases. The coefficient ratios of the quotient
for \(w_q\) decrease, so the series comparison rule gives
\(\Delta w\le0\). We also have \(\Delta j\le0\), completing
\eqref{eq:diff-comparisons} and the proof of \eqref{eq:diff-bound}.
If \(a\ne b\), its lower bound is positive. If \(a=b=c\), then
\(U=H_s(c)\) and \(U^{(4)}=c^2 U''\).
\end{proof}

\begin{proof}[Proof of Theorem~\ref{thm:product}]
For \(n\ge3\), Lemmas~\ref{lem:first} and \ref{lem:integral} and
Proposition~\ref{prop:differential} give \(B_j\ge0\) for all \(j\).
The integrand's differential factor at \(s=1\) is nonnegative by continuity.
Lemma~\ref{lem:remainder} and \eqref{eq:squares} now prove \eqref{eq:minorant}.
If \(x\ne y\) are positive, then \(B_1,B_2>0\).
Since \(S_n(z)-S_{n-1}(z)=z^{n-1}\), these two polynomials have no common
zero on \(\T\), and the minorant inequality is strict there. The minimum
principle gives strictness throughout the closed disk as well.

If \(x=y=e^{-a}\), then \(U_m=H_m(a)\), \(c=a\), and
\[
 L_m=m(\cosh a-1)\sinh a,\qquad E_n=0 .
\]
The first identity follows from
\(\sinh((m+1)a)+\sinh((m-1)a)=2\cosh a\,\sinh(ma)\);
the second follows from the difference \(H_{n+1}-H_{n-1}\) in
Lemma~\ref{lem:remainder}. Hence \(B_j=0\) for \(j\ge2\), and the right
side of \eqref{eq:squares} is \(B_1|S_n(z)|^2\). This gives exact contact at
all nontrivial \(n\)th roots of unity. Formula \eqref{eq:weight} extends
the bound and the equal-radius contact to zero radii by continuity.
The arithmetic mean dominates the geometric mean and \(P>0\).
Proposition~\ref{prop:flow} therefore applies with \(w=w_n\), \(\eta=1\),
and \(v_k=\sqrt2\,p_k\in\R\). Section~1 covers the two smaller orders.
\end{proof}

\end{document}
